\documentclass[10pt,a4paper]{amsart}
\usepackage[margin=19mm]{geometry}
\usepackage{amsmath,amssymb,mathtools}
\usepackage[scr=rsfs,cal=euler]{mathalfa}
\usepackage{xfrac}
\usepackage[unicode,hidelinks,bookmarks=false]{hyperref}
\newcommand{\ie}{i.e.\ }
\newcommand{\cf}{cf.\ }
\newcommand{\resp}{respectively\ }
\newcommand{\Subsection}{\subsection}
\providecommand{\nobreakdash}{\nobreak\hbox{-}}
\setlength{\emergencystretch}{2em}
\multlinegap0pt
\usepackage{bm}
\usepackage{bbm}
\usepackage{dsfont}
\usepackage{xspace}
\usepackage{cancel}
\usepackage{mathtools}
\usepackage{cleveref}
\usepackage[shortlabels]{enumitem}
\usepackage{amsthm}
\makeatletter
\@ifundefined{theorem}{\newtheorem{theorem}{Theorem}}{}
\@ifundefined{corollary}{\newtheorem{corollary}[theorem]{Corollary}}{}
\@ifundefined{lemma}{\newtheorem{lemma}[theorem]{Lemma}}{}
\@ifundefined{proposition}{\newtheorem{proposition}[theorem]{Proposition}}{}
\makeatother
\newcommand{\Ass}{\operatorname{Ass}}
\newcommand{\Ext}{\operatorname{Ext}}
\newcommand{\Frac}{\operatorname{Frac}}
\newcommand{\Hom}{\operatorname{Hom}}
\newcommand{\Jac}{\operatorname{Jac}}
\newcommand{\Spec}{\operatorname{Spec}}
\newcommand{\depth}{\operatorname{depth}}
\title[The Noetherian Case of Bayart's Power-Series Question]{The Noetherian Case of Bayart's Power-Series Question}
\author{Viet-Hoang Tran, Dung V. Nguyen, Quang X. Nguyen, Thieu N. Vo, Tan M. Nguyen}
\date{Source: arXiv:2608.12642v1 (CC BY 4.0)}
\begin{document}
\maketitle

\noindent\emph{Source and licence.} The Noetherian Case of Bayart's Power-Series Question. Version arXiv:2608.12642v1 (CC BY 4.0), distributed under the Creative Commons Attribution 4.0 International licence, as recorded in the arXiv licence metadata for this exact version and shown on the abstract page https://arxiv.org/abs/2608.12642v1. Full rights notice: licenses/arxiv-2608.12642-CC-BY-4.0.txt.\par\smallskip

\noindent\textit{Editorial scope.} Counted unit: the Corollary whose source label is \texttt{cor:reflexive-B-free} in the article (source0.tex lines 1047--1050, source label \texttt{cor:reflexive-B-free}), delivered together with the article's own complete original proof of it (source0.tex lines 1052--1083, one \texttt{proof} environment of 32 lines). No mathematics is written, repaired or re-derived: statement and proof are the article's own text line for line. Required context retained uncounted: lem:lifting-dual (lines 313--335); lem:movable-section (lines 719--729); prop:section-freeness (lines 808--826). Every definition, display and label of those lines is retained unchanged. Omitted from this package and not redistributed: 1--312, 336--718, 730--807, 827--1046, 1051--1051, 1084--1109. Nothing inside a retained excerpt is omitted or altered: every retained body is delivered exactly as the article's lines are written, so no omission receipt of any kind accompanies this package. Citations: the retained excerpts carry no citation command at all, so no bibliography entry of the article is transcribed and no reference is redistributed. Reference adaptation: none was needed and none was made; every reference inside the delivered unit resolves inside it, so no unresolved reference or citation is produced. Printed slips kept literal: no printed word, symbol, hypothesis or number of the counted statement or of its proof was altered. Layout and class: the article's own class is \texttt{amsart} with packages including amsmath, amssymb, amsfonts, geometry, fontenc, microtype, amsthm, mathtools, enumitem, xcolor, url, hyperref, cleveref; the wrapper is the lane's pinned \texttt{amsart} editorial wrapper at 10pt on A4, which restates the theorem-like environments the retained text needs and adds the article's own macro definitions that the retained text needs (\texttt{\textbackslash Ass}, \texttt{\textbackslash Ext}, \texttt{\textbackslash Frac}, \texttt{\textbackslash Hom}, \texttt{\textbackslash Jac}, \texttt{\textbackslash Spec}, \texttt{\textbackslash depth}), with extra wrapper packages (bm, bbm, dsfont, xspace, cancel, mathtools, cleveref, enumitem). The delivered numbering is the unit's own. Assets: no retained span contains a figure environment, a graphics-inclusion command, a PDF-page inclusion, a table or a TikZ picture; no image file is copied into this package, the package declares no asset dependency and no image file is redistributed with it. Licence: version arXiv:2608.12642v1 of this article is distributed under the Creative Commons Attribution 4.0 International licence, as recorded in the arXiv licence metadata for this exact version and shown on the abstract page https://arxiv.org/abs/2608.12642v1. A full rights notice is delivered at licenses/arxiv-2608.12642-CC-BY-4.0.txt. Characters: every retained excerpt is pure ASCII, so the delivered engine, XeLaTeX with its default Latin Modern text font, reports no missing character.
\begin{lemma}[Lifting freeness of a dual]
\label{lem:lifting-dual}
Let $S$ be a Noetherian ring, let $u\in\Jac(S)$ be a nonzerodivisor on
$S$, and suppose that
\[
  \overline S=S/uS
\]
is a domain.  Let $M$ be a finite $S$-module such that $u$ is a
nonzerodivisor on $M$, and put
\[
  \overline M=M/uM.
\]
Assume the following:
\begin{enumerate}[label=\textup{(\roman*)}]
  \item $\overline M$ is torsion-free over $\overline S$;
  \item $\overline M^\vee$ is a finite free $\overline S$-module;
  \item $\overline M_{\mathfrak p}$ is free over
  $\overline S_{\mathfrak p}$ for every
  $\mathfrak p\in\Spec(\overline S)$ satisfying
  $\depth\overline S_{\mathfrak p}\leq 2$.
\end{enumerate}
Then $M^\vee$ is a finite free $S$-module.
\end{lemma}

\begin{lemma}[A movable prime section]
\label{lem:movable-section}
For every finite $B$-module $T$, there is an integer $N\geq1$ such that,
with $z=y-x^N$, the following hold:
\begin{enumerate}[label=\textup{(\roman*)}]
  \item $z\in\Jac(B)$;
  \item $z$ is a nonzero prime element and $B/zB\cong A$;
  \item $z\notin\mathfrak q$ for every
        $\mathfrak q\in\Ass_B(T)\setminus V(x,y)$.
\end{enumerate}
\end{lemma}

\begin{proposition}[Low-depth freeness on a movable section]
\label{prop:section-freeness}
Let $F$ be a finite rank-one reflexive $B$-module, put
\[
  \Delta=\Ext^1_B(F^\vee,B),
\]
choose $z=y-x^N$ as in \cref{lem:movable-section} for $T=\Delta$, and set
$E=F/zF$.  Then:
\begin{enumerate}[label=\textup{(\roman*)}]
  \item $E$ is finite, torsion-free, and of rank one, and
        $E^\vee\cong A$;
  \item $E_x\cong A_x$;
  \item if $\mathfrak p\in\Spec(A)$ contains $x$ and
        $\depth A_{\mathfrak p}\leq2$, then
        $E_{\mathfrak p}\cong A_{\mathfrak p}$.
\end{enumerate}
Consequently, $E_{\mathfrak p}$ is free for every
$\mathfrak p\in\Spec(A)$ with $\depth A_{\mathfrak p}\leq2$.
\end{proposition}

\begin{corollary}[Reflexive modules over $B$]
\label{cor:reflexive-B-free}
Every finite rank-one reflexive $B$-module is free of rank one.
\end{corollary}

\begin{proof}
Let $F$ be such a module and put
\[
  \Delta=\Ext^1_B(F^\vee,B).
\]
As in the proof of \cref{prop:section-freeness}, the module $\Delta$ is
finite.  Choose $z=y-x^N$ by \cref{lem:movable-section} for
$T=\Delta$, and put $E=F/zF$.  By
\cref{prop:section-freeness}, the module $E$ is torsion-free,
$E^\vee\cong A$, and $E_{\mathfrak p}$ is free whenever
$\depth A_{\mathfrak p}\leq2$.

The other hypotheses of \cref{lem:lifting-dual} follow from
\cref{lem:movable-section}: the element $z$ belongs to $\Jac(B)$, is
regular on the domain $B$, and satisfies $B/zB\cong A$; it is also regular
on the torsion-free module $F$.  Apply \cref{lem:lifting-dual} with
\[
 S=B,\qquad u=z,\qquad M=F,
 \qquad\overline S=A,\qquad\overline M=E.
\]
The lemma shows that $F^\vee$ is free.  Localization at the generic point gives
\[
 F^\vee\otimes_B\Frac(B)
 \cong\Hom_{\Frac(B)}
       (F\otimes_B\Frac(B),\Frac(B))
 \cong\Frac(B),
\]
so $F^\vee$ has rank one and is isomorphic to $B$.  Reflexivity now gives
\[
  F\cong F^{\vee\vee}\cong B.
\]
\end{proof}

\end{document}
